\begin{lemma}[Mehler averaging formula and positivity preservation]
Let $s > 0$ and $v \in L^2(\gamma_n)$, and write $r = e^{-s} \in (0,1)$ and
$\tau = \sqrt{1-e^{-2s}} = \sqrt{1-r^2} \in (0,1]$.  Then the spectrally defined
semigroup $P_s$ of \noderef{OUSemigroup} is given by the probabilistic averaging
formula
\[
  P_s v(x)
    = \int_{\mathbb{R}^n} v\bigl(r\,x + \tau\,y\bigr)\, d\gamma_n(y)
  \qquad\text{for } \gamma_n\text{-a.e. } x \in \mathbb{R}^n ,
\]
the inner integral being absolutely convergent for $\gamma_n$-a.e.\ $x$.
Consequently $P_s$ is positivity preserving: if $v \geq 0$ $\gamma_n$-a.e.\ then
$P_s v \geq 0$ $\gamma_n$-a.e.
\end{lemma}
\begin{proof}
Define, for $v \in L^2(\gamma_n)$,
\[
  (Q_s v)(x) = \int_{\mathbb{R}^n} v(r x + \tau y)\, d\gamma_n(y) .
\]

\emph{Step 1: the affine map pushes $\gamma_n\otimes\gamma_n$ to a measure with second marginal $\gamma_n$.}
By \noderef{GaussMeasure} the density of $\gamma_n$ factorizes over the
coordinates as $\prod_j \pi^{-1/2}e^{-x_j^2}$, and $\gamma_n$ is a probability
measure by \noderef{GaussIsProb}; thus $\gamma_n$ is the $n$-fold product of the
one-dimensional measure $d\gamma_1(t) = \pi^{-1/2}e^{-t^2}dt$, which is the
centred Gaussian law of variance $1/2$ on $\mathbb{R}$.  Consequently, if $X$ and
$Y$ are independent with law $\gamma_n$, then each coordinate of
$rX + \tau Y$ is a sum of two independent centred Gaussians of variances
$r^2/2$ and $\tau^2/2$, hence centred Gaussian of variance
$(r^2+\tau^2)/2 = 1/2$, and distinct coordinates remain independent.  Therefore
\begin{equation*}
  rX + \tau Y \ \text{ has law } \ \gamma_n . \tag{1}
\end{equation*}

\emph{Step 2: $Q_s$ is a well-defined contraction on $L^2(\gamma_n)$.}
By Tonelli and (1),
\[
  \iint |v(rx+\tau y)|^2\, d\gamma_n(y)\, d\gamma_n(x)
    = \int |v|^2\, d\gamma_n = \|v\|_2^2 < \infty .
\]
Hence for $\gamma_n$-a.e.\ $x$ the function $y \mapsto v(rx+\tau y)$ is in
$L^2(\gamma_n)$, so in $L^1(\gamma_n)$ by \noderef{L2IntoL1}; thus $Q_sv(x)$ is
defined by an absolutely convergent integral for a.e.\ $x$.  Jensen's inequality
for the probability measure $\gamma_n$ gives
$|Q_sv(x)|^2 \leq \int |v(rx+\tau y)|^2 d\gamma_n(y)$, and integrating in $x$ and
using the display above,
\begin{equation*}
  \|Q_s v\|_{L^2(\gamma_n)} \leq \|v\|_{L^2(\gamma_n)} . \tag{2}
\end{equation*}
In particular $Q_s$ is a bounded linear operator on $L^2(\gamma_n)$.  The
argument of this step used nothing about $v$ beyond $v \in L^2(\gamma_n)$, so
its conclusions hold with any element of $L^2(\gamma_n)$ in place of $v$.

\emph{Step 3: $Q_s$ and $P_s$ agree on each Hermite function.}
Fix $\alpha \in \mathbb{N}_0^n$.  By \noderef{HermiteFun},
$h_\alpha(z) = \prod_{j} h_{\alpha_j}(z_j)$, and the $j$th coordinate of
$rx+\tau y$ is $r x_j + \tau y_j$.  Hence
\[
  h_\alpha(rx+\tau y) = \prod_{j=1}^n h_{\alpha_j}(r x_j + \tau y_j),
\]
a product of functions of the separate coordinates $y_j$.

We justify the factorization of the integral.  By \noderef{HermiteFunMemLp},
$h_\alpha \in L^2(\gamma_n)$, so Step 2 applied with $v = h_\alpha$ shows that for
$\gamma_n$-a.e.\ $x$ the function $y \mapsto h_\alpha(rx+\tau y)$ is
$\gamma_n$-integrable; fix such an $x$.  Tonelli's theorem for the product measure
$\gamma_n = \gamma_1^{\otimes n}$ of Step 1, applied to the nonnegative function
$y \mapsto \prod_j |h_{\alpha_j}(rx_j+\tau y_j)|$, gives
\[
  \int \prod_{j=1}^n \bigl|h_{\alpha_j}(rx_j+\tau y_j)\bigr|\, d\gamma_n(y)
    = \prod_{j=1}^n \int_{\mathbb{R}} \bigl|h_{\alpha_j}(rx_j+\tau\sigma)\bigr|\,
        d\gamma_1(\sigma)
\]
as an identity in $[0,+\infty]$, and the left-hand side is finite by the previous
sentence.  No factor on the right vanishes: by \noderef{Hermite1} and the
recurrence recorded there, $\mathrm{He}_{\alpha_j}$ is monic of degree $\alpha_j$, so
$h_{\alpha_j}$ is a nonzero polynomial; since $\tau > 0$, the map
$\sigma \mapsto h_{\alpha_j}(rx_j+\tau\sigma)$ is again a nonzero polynomial, hence
nonzero off a finite set, and $\gamma_1$ has a strictly positive density by
\noderef{GaussMeasure}, so its integral against $\gamma_1$ is strictly positive.
A product of elements of $[0,+\infty]$ with no vanishing factor is finite only if
every factor is finite, so each one-dimensional integral above is finite.
Fubini's theorem for the product measure therefore gives
\[
  (Q_s h_\alpha)(x)
    = \prod_{j=1}^n \int_{\mathbb{R}} \pi^{-1/2}e^{-\sigma^2}\,
        h_{\alpha_j}(r x_j + \tau \sigma)\, d\sigma .
\]
Since $0 < r < 1$, \noderef{Hermite1Reproduction} applies with this $r$ (note
$\tau = \sqrt{1-r^2}$ as required there) and evaluates the $j$th factor as
$r^{\alpha_j} h_{\alpha_j}(x_j)$.  Multiplying,
\[
  (Q_s h_\alpha)(x)
    = r^{\sum_j \alpha_j}\prod_j h_{\alpha_j}(x_j)
    = r^{|\alpha|}\, h_\alpha(x)
    = e^{-s|\alpha|}\, h_\alpha(x),
\]
using \noderef{MultiIndexSum} and $r = e^{-s}$.  On the other hand
\noderef{OUSemigroupSpectral} (applicable since $s > 0 \geq 0$) gives
$P_s h_\alpha = e^{-s|\alpha|}h_\alpha$, because the coefficients of $h_\beta$ are
$\delta_{\alpha\beta}$ by \noderef{HermiteL2Orthonormal}.  Hence
$Q_s h_\alpha = P_s h_\alpha$ for every $\alpha$.

\emph{Step 4: $Q_s = P_s$ on all of $L^2(\gamma_n)$.}
Both operators are linear and bounded on $L^2(\gamma_n)$: $Q_s$ by (2), and
$P_s$ by \noderef{OUSelfAdjointContraction}, which gives
$\|P_s v\|_2 \leq \|v\|_2$.  They agree on every $h_\alpha$ by Step 3, hence on
the linear span of $\{h_\alpha\}$ by linearity, hence on the closure of that span
by continuity.  By \noderef{HermiteL2Complete} that closure is all of
$L^2(\gamma_n)$.  Therefore $Q_s v = P_s v$ in $L^2(\gamma_n)$, i.e.\
$P_s v(x) = \int v(rx+\tau y)\, d\gamma_n(y)$ for $\gamma_n$-a.e.\ $x$, which is
the asserted formula.

\emph{Step 5: positivity preservation.}
Suppose $v \geq 0$ $\gamma_n$-a.e., and let $N$ be the $\gamma_n$-null set where
this fails.  By (1) the map $(x,y) \mapsto rx+\tau y$ pushes
$\gamma_n\otimes\gamma_n$ forward to $\gamma_n$, so the set of pairs $(x,y)$ with
$rx+\tau y \in N$ has $(\gamma_n\otimes\gamma_n)$-measure $\gamma_n(N) = 0$.  By
Fubini for this null set, for $\gamma_n$-a.e.\ $x$ we have $rx+\tau y \notin N$
for $\gamma_n$-a.e.\ $y$, and then $v(rx+\tau y) \geq 0$ for $\gamma_n$-a.e.\ $y$.
For such $x$ the integral defining $Q_sv(x)$ is an integral of an a.e.\
nonnegative integrable function, hence $Q_s v(x) \geq 0$.  Combining with Step 4,
$P_s v \geq 0$ $\gamma_n$-a.e.
\end{proof}
